\chapter{Conclusion}
\label{ch:conclusion}

\section{Summary}

This dissertation implemented a visual behaviour-tree plugin in Godot 4.6 and studied the display problem of large behaviour trees. Godot provides general graph-editing controls but does not provide a complete behaviour-tree editing and execution workflow. Traditional node graphs also occupy a large amount of canvas space as nodes, parameters and debugging information increase.

The final plugin can create, save, load and execute behaviour trees, and supports a blackboard, Decorators, multiple composite-node types, Actions, Conditions, undo, redo and Live Debug. A complex two-dimensional scene uses a 241-node behaviour tree to control enemy detection, defence, melee attack, ranged attack, jumping, climbing, search, retreat, healing and patrol.

Experiments across five tree scales showed that Compact reduced card area by 61.09--61.35\%, Overview reduced it by 90.27--90.34\%, and Search dimmed 96.15--99.67\% of non-target cards. Focus and Collapse also substantially reduced the context currently displayed, and no condition changed saved positions or execution order.

The three-screen experiment showed that complex-tree coverage on the 31.55-inch screen was 2.54 times that on the 15.94-inch screen. The effects of Compact, Overview and Search remained stable across different screens. Overall, the display optimisations improved some objective display and editing-interaction measures for large behaviour trees.
